\label{sec:pert}

For completeness, we collect the main steps of the original
derivations of \citeres{Luscher:2010iy,Luscher:2011bx} in this section.


\subsection{Definition of the gradient flow}

In the following, we work in $D$-dimensional Euclidean space-time
with $D=4-2\ep$.  The \gff\ continues the gluon and quark fields
$A^a_\mu(x)$ and $\psi_\alpha^i(x)$ of regular\footnote{We use the terms
  ``flowed'' and ``regular'' \qcd\ to distinguish quantities defined at
  $t>0$ from those defined at $t=0$.} \qcd\ to $(D+1)$-dimensional
fields $B^a_\mu(t,x)$ and $\chi_\alpha^i(t,x$) through the boundary
conditions
\begin{equation}
  \begin{split}
    B_\mu^a (t=0,x) = A_\mu^a (x)\,,\qquad \chi_\alpha^i (t=0,x)=
    \psi_\alpha^i(x)
    \label{eq:bound}
  \end{split}
\end{equation}
and the flow equations\,\cite{Luscher:2010iy,Luscher:2013cpa}
\begin{equation}
  \begin{split}
    \partial_t B_\mu^a &= \mathcal{D}^{ab}_\nu G_{\nu\mu}^b + \kappa
    \mathcal{D}^{ab}_\mu \partial_\nu B_\nu^b\,,\\ \partial_t \chi &= \Delta \chi
    - \kappa \partial_\mu B_\mu^a T^a \chi\,,\\ \partial_t \bar
    \chi &= \bar \chi \overleftarrow \Delta + \kappa \bar \chi
    \partial_\mu B_\mu^a T^a\,,
    \label{eq:flow}
  \end{split}
\end{equation}
where the ``flow time'' $t$ is a parameter of mass dimension minus two,
and $\kappa$ is a gauge parameter which drops out of physical
observables (see below).

The $(D+1)$-dimensional field-strength tensor is defined as
\begin{align}
  G_{\mu\nu}^a = \partial_\mu B_\nu^a - \partial_\nu B_\mu^a + f^{abc}
  B_\mu^b B_\nu^c\,,
\end{align}
the covariant derivative in the adjoint representation is given by
\begin{align}
  \mathcal{D}_\mu^{ab} = \delta^{ab} \partial_\mu - f^{abc} B_\mu^c\,,
\end{align}
and
\begin{equation}
  \begin{split}
    \Delta = \mathcal{D}^\mathrm{F}_\mu \mathcal{D}^\mathrm{F}_\mu
    \,,\qquad
    \overleftarrow{\Delta} = \overleftarrow{\mathcal{D}}^\mathrm{F}_\mu
    \overleftarrow{\mathcal{D}}^\mathrm{F}_\mu \,,
  \end{split}
\end{equation}
with the covariant derivative in the fundamental representation,
\begin{equation}
  \mathcal{D}^\mathrm{F}_\mu = \partial_\mu + B^a_\mu T^a \,,\qquad
  \overleftarrow{\mathcal{D}}^\mathrm{F}_\mu =
    \overleftarrow{\partial}_\mu - B^a_\mu T^a \,.
\end{equation}
As usual, the color indices of the adjoint representation are denoted by
$a,b,c,\ldots$, while $\mu,\nu,\rho,\ldots$ are $D$-dimensional Lorentz
indices. Color indices of the fundamental representation are denoted by
$i,j,k,\ldots$, but they are suppressed throughout this paper, unless
required by clarity; similarly for spinor indices
$\alpha,\beta,\gamma,\ldots$. The symmetry generators $T^a$ are understood in
the fundamental representation. The obey the commutation relation
\begin{equation}
  \begin{split}
    [T^a,T^b] = f^{abc}T^c\,,
  \end{split}
\end{equation}
with the structure constants $f^{abc}$ and the trace is normalized to
\begin{equation}
  \mathrm{Tr}(T^a T^b) = - \TR \delta^{ab} \, ,
\end{equation}
with $\TR > 0$.

Different choices of $\kappa$ in \eqn{eq:flow} correspond to gauge
transformations of the form
\begin{align}
  \chi \rightarrow \chi^\prime = \Lambda \chi \quad \quad \text{and}
  \quad \quad B_\mu^a T^a \rightarrow B_\mu^{a\prime} T^a= \Lambda
  B_\mu^a T^a \Lambda^{-1} + \Lambda \partial_\mu \Lambda^{-1}\,,
\end{align}
where
\begin{align}
  \Lambda(t,x)=e^{-\int_0^t \mathrm{d}s \, \kappa \partial_\mu B_\mu^a(s,x) T^a}\,.
\end{align}
Of course, all observables are independent of the gauge parameter
$\kappa$\,\cite{Luscher:2010iy}.  In perturbative calculations, it is
usually most convenient to set $\kappa = 1$.

The flow equations \noeqn{eq:flow} can be incorporated in a Lagrangian
formalism by defining
\begin{align}
  \mathcal{L} = \mathcal{L}_\mathrm{\qcd} +
  \mathcal{L}_\textrm{gauge-fixing} + \mathcal{L}_\mathrm{ghost} +
  \mathcal{L}_{B} + \mathcal{L}_{\chi}.
  \label{eq:Lagrangian_complete}
\end{align}
The first three terms constitute the regular Yang-Mills Lagrangian, with
fermions added in the fundamental representation (quarks). Introducing an
index $f$ in order to distinguish different quark flavors of mass $m_f$,
the classical, gauge-fixing, and Faddeev-Popov ghost part are given by
\begin{equation}
  \begin{split}
  \mathcal{L}_\mathrm{\qcd} &= \frac{1}{4g^2} F_{\mu \nu}^a F_{\mu
    \nu}^a + \sum_{f=1}^{\NF}\bar{\psi}_f ( \slashed{D}^\mathrm{F}+m_f)
  \psi_f\,,\\ \mathcal{L}_\textrm{gauge-fixing}&= \frac{1}{2 g^2 \xi}
  (\partial_\mu A_\mu^a)^2\,,\\ \mathcal{L}_\mathrm{ghost} &=
  \frac{1}{g^2}\partial_\mu \bar{c}^a D_\mu^{ab} c^b\,,
  \label{eq:lags}
  \end{split}
\end{equation}
respectively, where
\begin{equation}
  \begin{split}
    F_{\mu\nu}^a = \partial_\mu A^a_\nu - \partial_\nu A^a_\mu + f^{abc}A_\mu^bA_\nu^c
    %% \label{eq:}
  \end{split}
\end{equation}
is the regular field strength tensor and
\begin{equation}
  D_\mu^\mathrm{F} = \partial_\mu + A_\mu^a T^a \,,\qquad D_\mu^{ab} = \delta^{ab} \partial_\mu - f^{abc} A_\mu^c
\end{equation}
are the regular covariant derivatives in the fundamental and adjoint representation, respectively, $g$
is the gauge coupling, $\xi$ the \qcd\ gauge parameter, and $\NF$ the
number of different quark flavors.  The flow equations are incorporated
by introducing Lagrange multiplier fields
\begin{equation}
  \begin{split}
    L_\mu^a(t,x)\qquad\text{and}\qquad
    \lambda_f(t,x)\,,\ \bar\lambda_f(t,x)\,,
  \end{split}
\end{equation}
of mass dimensions 3 and 5/2 that otherwise carry the same quantum
numbers as the flowed gluon and quark/antiquark
fields $B_\mu^a$ and $\chi,\bar\chi$, respectively.
Their Euler-Lagrange equations derived from
\begin{equation}
  \begin{split}
  \mathcal{L}_{B} &= -2 \int_0^\infty \mathrm{d}t \, \textrm{Tr} \left[
    L_\mu^a T^a \left(\partial_t B_\mu^b T^b -\mathcal{D}_\nu^{bc}
    G_{\nu \mu}^c T^b - \kappa \mathcal{D}_\mu^{bc} \partial_\nu B_\nu^c
    T^b \right) \right]\,,\\ \mathcal{L}_{\chi} & =
  \sum_{f=1}^{\NF}\int_0^\infty \mathrm{d}t \, \Big( \bar{\lambda}_f
  \left(\partial_t - \Delta + \kappa \left(\partial_\mu B_\mu^a\right)
  T^a \right) \chi_f \\&\hspace*{10em}
  + \bar{\chi}_f \left(\overleftarrow{\partial_t} -
  \overleftarrow{\Delta} - \kappa \left(\partial_\mu B_\mu^a\right) T^a
  \right) \lambda_f \Big)\,,
  \label{eq:Lagrangian_flow}
  \end{split}
\end{equation}
indeed lead to \eqn{eq:flow}\,\cite{Luscher:2011bx}.

Note that the Lagrangian \noeqn{eq:Lagrangian_complete} does not include
flowed ghost fields $d^a(t,x)$ and $\bar{d}^a(t,x)$. They arise in the
same way as the usual Faddeev-Popov ghosts $c^a(x)$ and $\bar{c}^a(x)$
due to gauge-fixing, and obey the initial condition
\begin{align}
  \left. d^a (t,x) \right|_{t=0} = c^a(x).
\end{align}
Similar to the ghosts of the regular gauge fields, they always form
closed loops as long as there are no external $d^a$ or $\bar{d}^a$
fields (which can be avoided by considering only physical degrees of
freedom in amplitudes with external gluons). As becomes clear later,
closed loops of only flowed fields vanish, so that one can omit $d^a$
and $\bar{d}^a$ already at the level of the Lagrangian.


\subsection{Perturbative solution of the flow equations}

Let us introduce the short-hand notation
\begin{equation}
  \begin{split}
    \int_p \equiv \int\frac{\dd^D p}{(2\pi)^D}\,,\qquad
    \int_x \equiv \int\dd^D x\,,
    %% \label{eq:}
  \end{split}
\end{equation}
where it should be clear from the context whether an integration
variable is in position ($x,y,z,\ldots$) or momentum space
($p,k,q,\ldots$).

It is helpful to separate the flow equation \noeqn{eq:flow} of the flowed gauge
field $B_\mu^a(t,x)$ into a linear and a non-linear part:
\begin{equation}
  \begin{split}
  &\partial_t B_\mu^a = \partial_\nu \partial_\nu B_\mu^a + (\kappa - 1)
  \partial_\mu \partial_\nu B_\nu^a + R_\mu^a, \\ &R_\mu^a = 2 f^{abc}
  B_\nu^b \partial_\nu B_\mu^c - f^{abc} B_\nu^b \partial_\mu B_\nu^c +
  (\kappa - 1) f^{abc} B_\mu^b \partial_\nu B_\nu^c + f^{abe} f^{cde}
  B_\nu^b B_\nu^c B_\mu^d.
  \label{eq:flowsep}
  \end{split}
\end{equation}
The linear equation can be solved by introducing the integration kernel
\begin{equation}
  K_{\mu \nu} (t,x) = \int_p \frac{e^{\mathrm{i}px}}{p^2} \left( \left(
  \delta_{\mu \nu} p^2 - p_\mu p_\nu \right) e^{-tp^2} + p_\mu p_\nu \,
  e^{-\kappa tp^2} \right) \equiv \int_p e^{\ii px}\widetilde K_{\mu\nu}(t,p)\,,
  \label{eq:kmunu}
\end{equation}
which fulfills
\begin{equation}
  \lim_{t \to 0} K_{\mu\nu}(t,x) = \delta_{\mu\nu} \delta^{(D)}(x).
\end{equation}
Taking into account the initial condition \noeqn{eq:bound}, the full
solution of the flow equation is then given by
\begin{equation}
  B_\mu^a (t,x) = \int_y K_{\mu \nu}(t,x-y) A_{\nu}^a(y)+ \int_y
  \int_0^t \mathrm{d}s \, K_{\mu \nu}(t-s,x-y) R_\nu^a(s,y)\,,
\end{equation}
or, in momentum space,
\begin{equation}
  \widetilde{B}_\mu^a(t,p)= \int_x e^{-\mathrm{i}px} B_\mu^a(t,x)=
  \widetilde{K}_{\mu \nu}(t,p) \widetilde{A}_\nu^a(p)+ \int_0^t
  \mathrm{d}s \, \widetilde{K}_{\mu \nu}(t-s,p)
  \widetilde{R}_\nu^a(s,p).
  \label{eq:gluon_solution_momentum}
\end{equation}
By inserting the solution iteratively into itself, one can express the
Fourier transform of the non-linear part of \eqn{eq:flowsep} as
\begin{equation}
  \begin{split}
  \widetilde{R}_\mu^a(t,p) &=
  \int_{q,l,k} (2 \pi)^D \, \delta^{(D)}(p-q-l-k)\bigg[
   \delta^{(D)}(k)\cdot X_{2,\mu\nu\rho}^{abc}(q,l)\, \widetilde{B}_{\nu}^{b}(t,q)
      \widetilde{B}_{\rho}^{c}(t,l)\\&
  +
   X_{3,\mu\nu\rho\sigma}^{abcd}\, \widetilde{B}_{\nu}^{b}(t,q)
      \widetilde{B}_{\rho}^{c}(t,l)
      \widetilde{B}_{\sigma}^{d}(t,k)\bigg]\,,
    %% \label{eq:}
  \end{split}
\end{equation}
where
\begin{equation}
  \begin{split}
  X_{2,\mu \nu \rho}^{abc}(q,l) &= -\mathrm{i} f^{a b c} \big[
  (l-q)_\mu \delta_{\nu \rho} + 2q{}_{\rho} \delta_{\mu \nu}- 2l{}_{\nu}
  \delta_{\mu \rho} + (\kappa - 1) ( q{}_{\nu} \delta_{\mu
    \rho} -l{}_{\rho} \delta_{\mu \nu}) \big]\,,\\
  X_{3,\mu \nu \rho \sigma}^{abcd} &= f^{abe}
  f^{cde}(\delta_{\mu \sigma} \delta_{\nu \rho} - \delta_{\mu \rho}
  \delta_{\nu \sigma}) + f^{ade} f^{bce}(\delta_{\mu \rho}
  \delta_{\nu \sigma}- \delta_{\mu \nu} \delta_{\rho \sigma})\\&\quad
  + f^{ace} f^{dbe}(\delta_{\mu \nu} \delta_{\rho \sigma}-
  \delta_{\mu \sigma} \delta_{\nu \rho}).
    %% \label{eq:}
  \end{split}
\end{equation}
The structure of $X_3$ is identical to the four-gluon vertex of regular
\qcd. When formulating the Feynman rules later, $X_2$ and $X_3$ describe
the three- and four-point vertices of the flowed gluon fields.

The flow equation \noeqn{eq:flow} for the flowed quark fields can be
solved by again splitting it into a linear and a non-linear part,
\begin{align}
  \partial_t \chi =\partial_\mu \partial_\mu \chi + \Delta^\prime \chi
  \quad \textrm{with} \quad \Delta^\prime = (1-\kappa) \partial_\mu
  B_\mu^a T^a + 2 B_\mu^a T^a \partial_\mu + B_\mu^a B_\mu^b T^a T^b\,.
\end{align}
The linear equation is solved by the integration kernel
\begin{equation}
  K(t,x)=\int_p e^{\mathrm{i}px} e^{-tp^2} \equiv \int_p e^{\ii
    px}\widetilde K(t,p)\,,
  \label{eq:k}
\end{equation}
with the help of which we can write the full solution as
\begin{equation}
  \chi(t,x)= \int_y K(t,x-y)\psi(y) + \int_y \int_0^t \mathrm{d}s \, K(t-s,x-y) \Delta^\prime \chi(s,y).
\end{equation}
Here and in what follows, we suppress the flavor index $f$ unless
required for clarity.  The non-linear part of the Fourier-transformed
field
\begin{equation}
  \widetilde{\chi}(t,p)= \widetilde{K}(t,p)\widetilde{\psi}(p) + \int_0^t \mathrm{d}s \, \widetilde{K}(t-s,p)\widetilde{\Delta^\prime \chi}(s,p)
  \label{eq:quark_solution_momentum}
\end{equation}
can be expressed as
\begin{equation}
  \begin{split}
    \widetilde{\Delta^\prime \chi}(t,p) =&
    \int_{q,r} (2\pi)^D \delta^{(D)}(p-q-l-r)\bigg[
      \delta^{(D)}(l)\cdot Y_{1,\nu}^{b}(p,q,r)
      \widetilde{B}_{\nu}^{b}(t,q)
      \\&+
    \frac{1}{2}
    Y_{2,\nu\rho}^{bc}(p,q,l,r)
    \widetilde{B}_{\nu}^{b}(t,q)
    \widetilde{B}_{\rho}^{c}(t,l)\bigg] \widetilde{\chi}(t,r)\,,
  \end{split}
\end{equation}
where
\begin{equation}
  \begin{split}
  Y_{1,\nu}^b(q,r) &= \mathrm{i} \big( 2 r_\nu + (1-\kappa) q_\nu \big)
  T^b\,,\qquad
  Y_{2,\nu\rho}^{bc} = \delta_{\nu\rho} \bigl\{ T^{b}, T^{c} \bigr\}\,.
  \end{split}
\end{equation}
These expressions lead to the three- and four-point vertices of
the flowed quark fields. For $\bar{\chi}$ one proceeds analogously.


\subsection{Feynman rules}
\label{ssec:Feynman_rules}


\subsubsection*{(Flowed) Propagators}

Plugging in the solution of the flowed gluon field in momentum space
\eqn{eq:gluon_solution_momentum} into the two-point function, one
finds\,\cite{Luscher:2011bx}
\begin{align}
  \left\langle \widetilde{B}_\mu^a(t,p) \widetilde{B}_\nu^b(s,q)
  \right\rangle \Big|_{\textrm{LO}} &= \widetilde{K}_{\mu \rho}(t,p)
  \widetilde{K}_{\nu \sigma}(s,q) \left\langle \widetilde{A}_\rho^a(p)
  \widetilde{A}_\sigma^b(q) \right\rangle \nonumber \\ &= (2\pi)^D
  \delta^{(D)}(p+q) \, g^2\,D_{\mu\nu}^{ab}(p,t+s,\xi,\kappa)\,,
\end{align}
where
\begin{equation}
  \begin{split}
    D_{\mu\nu}^{ab}(p,t,\xi,\kappa) =
  \delta^{ab} \frac{1}{p^2} \bigg( \Big(
  \delta_{\mu \nu} - \frac{p_\mu p_\nu}{p^2} \Big) e^{-tp^2} + \xi
  \frac{p_\mu p_\nu}{p^2} e^{-\kappa tp^2} \bigg)\,,
    %% \label{eq:}
  \end{split}
\end{equation}
and we have used the result for the fundamental gluon propagator,
\begin{equation}
\left\langle \widetilde{A}_\mu^a(p) \widetilde{A}_\nu^b(q) \right\rangle
\Big|_{\text{LO}} = (2\pi)^D \delta^{(D)}(p+q) \, g^2
D_{\mu\nu}^{ab}(p,0,\xi,0)\,.
\end{equation}
The factor $g^2$ is taken into account in the corresponding
vertices further below. Since the flowed gluon propagator coincides with the
fundamental gluon propagator at $t+s=0$, we can express both of them by
the same Feynman rule:
\begin{align}
  & \raisebox{-0.5em}{\includegraphics{images/gluon-propagator.pdf}}
  & = &
  \begin{split}
    D_{\mu\nu}^{ab}(p,t+s,\xi,\kappa) \,.
  \end{split}
  & &
  \label{eq:flowedgluon}
\end{align}
We refer to this as the (flowed) gluon propagator. \eqn{eq:flowedgluon}
also applies to the mixed propagator $\langle \widetilde A \widetilde B
\rangle$, which is obtained by setting one of the two flow-time
variables to zero.

The same can be done for flowed quark fields. Inserting their momentum space
solution in \eqn{eq:quark_solution_momentum} into the two-point
function results in
\begin{align}
  \left. \left\langle \widetilde{\chi}^i_\alpha(t,p)
  \widetilde{\bar{\chi}}^j_\beta(s,q) \right\rangle
  \right|_{\text{LO}} &= \widetilde{K}(t,p) \widetilde{K}(s,q)
  \left\langle \widetilde{\psi}^i_\alpha(p)
  \widetilde{\bar{\psi}}^j_\beta(q) \right\rangle \nonumber \\ &=
  (2\pi)^D \delta^{(D)}(p+q)\, S^{ij}_{\text{F},\alpha\beta}(p,m,t+s)\,,
\end{align}
where
\begin{equation}
  \begin{split}
S^{ij}_\mathrm{F}(p,m,t) &= \delta^{ij} \frac{-i\slashed{p}+m}{p^2+m^2}\,
e^{-tp^2},
  \end{split}
\end{equation}
and we have used the result for the fundamental quark propagator
\begin{equation}
  \left\langle \widetilde{\psi}^i_\alpha(p)
  \widetilde{\bar{\psi}}^j_\beta(q) \right\rangle = (2\pi)^D
  \delta^{(D)}(p+q)\,S_{\text{F},\alpha\beta}^{ij}(p,m,0)\,.
\end{equation}
Since one can express the fundamental propagator through the flowed
propagator at vanishing flow times, both can be represented by the same
Feynman rule:
\begin{align}
  & \raisebox{-0.5em}{\includegraphics{images/quark-propagator.pdf}}
  & = &
  \begin{split}
    S_{\text{F},\alpha\beta}^{ij}(p,m,t+s) \,.
  \end{split}
  & &
\end{align}
We refer to this as the (flowed) quark propagator, which again
includes the mixed propagators
$\langle\psi\bar\chi\rangle$ and $\langle\chi\bar\psi\rangle$.


\subsubsection*{Flow lines}

Since there are no quadratic terms of the Lagrange multiplier fields
$L_\mu^a$, $\lambda$, and $\bar{\lambda}$ in \eqn{eq:Lagrangian_flow},
there are no propagators for these fields.  However, the
Lagrangian contains bilinear terms of a flowed field and a Lagrange
multiplier field.  Using standard methods, one derives the two-point
function at leading order as
\begin{equation}
  \left\langle B_\mu^a(t,x) L_\nu^b(s,y) \right\rangle = \delta^{ab}
  H_{\mu\nu}(t-s,x-y)\,,
\end{equation}
where $H_{\mu\nu}(t,x)$ obeys the equation
\begin{equation}
  \left[(\partial_t - \partial_\sigma \partial_\sigma) \delta_{\mu\nu} + (1 -
  \kappa) \partial_\mu \partial_\nu\right] H_{\nu\rho}(t,x) =
  \frac{1}{2\TR} \delta_{\mu\rho} \delta(t) \delta^{(D)}(x)\,,
\end{equation}
and $\TR$ is the usual color trace normalization.  Since the fundamental
gluon field $A_\mu^a(x) = B_\mu^a(0,x)$ does not couple to the Lagrange
multiplier field $L_\mu^a$, and all flow-time variables are positive,
one can impose the initial condition
\begin{equation}
  \begin{split}
  \left\langle B_\mu^a(t,x) L_\nu^b(s,y) \right\rangle
  \bigg|_{t=0} = 0 \quad \Rightarrow \quad
  H_{\mu\nu}(-s,x) = 0\,,
    \label{eq:bl}
  \end{split}
\end{equation}
so that the unique solution becomes
\begin{equation}
  H_{\mu\nu}(t,x) = \frac{1}{2\TR} \theta(t) K_{\mu\nu}(t,x)\,.
  \label{eq:hmunu}
\end{equation}
We refer to the $\langle BL\rangle$ bilinear as ``gluon flow line''.
The inverse of the factor $1/(2\TR)$ appears in the
corresponding ``flow vertices'' further below. We can therefore discard it
altogether and write
\begin{align}
  & \raisebox{-1.2em}{\includegraphics{images/gluon-flowline.pdf}}
  & = &
  \begin{split}
    \delta^{ab} \, \theta(t-s)\,\widetilde K_{\mu\nu}(t-s,p) \,,
  \end{split}
  & &
  \label{eq:blfr}
\end{align}
where $\widetilde K_{\mu\nu}(t,p)$ has been defined in \eqn{eq:kmunu},
and the adjacent arrow indicates the direction towards increasing
flow time as implied by the $\theta$-distribution.  As opposed to an
actual propagator, the gluon flow line is a regular function for all
$p$.

Similarly, one determines the mixed fermionic two-point function
at leading order as
\begin{equation}
  \left\langle \chi_\alpha^i(t,x) \bar{\lambda}_\beta^j(s,y) \right\rangle =
  \delta_{\alpha\beta} \delta^{ij} G(t-s,x-y)\,,
\end{equation}
where $G(t,x)$ obeys
\begin{equation}
  (\partial_t - \partial_\mu \partial_\mu) G(t,x) = \delta(t) \delta^{(D)}(x)\,,
\end{equation}
with the condition
\begin{equation}
  \left\langle \chi_\alpha^i(t,x) \bar{\lambda}_\beta^j(s,y)
  \right\rangle \bigg|_{t=0} = 0 \quad \Rightarrow \quad G(-s,x)= 0\,.
\end{equation}
The unique solution reads
\begin{equation}
  G(t,x) = \theta(t) K(t,x)\,.
\end{equation}
The Fourier transformed expression
defines the ``fermion flow line'' Feynman rule, where the
$\theta$-distribution again imposes a direction as indicated by the adjacent arrow
pointing towards increasing flow time:
\begin{align}
  & \raisebox{-1.2em}{\includegraphics{images/quark-flowline.pdf}}
  & = &
  \begin{split}
    \delta_{\alpha\beta} \delta_{ij} \, \theta(t-s) \, \widetilde K(t-s,p) \,,
  \end{split}
  & &
  \label{eq:chibarlambda}
\end{align}
where $\widetilde K(t,p)$ has been defined in \eqn{eq:k}.  As usual, the
arrow \textit{on} the fermion line denotes the ``charge flow'' of the
fermion. In this way, we have a unified Feynman rule for both the
$\langle\chi\bar\lambda\rangle$ and the $\langle\lambda\bar\chi\rangle$
bilinear, where the latter can be obtained analogously as above
and simply corresponds to reversing the direction of the charge flow:
\begin{align}
  & \raisebox{-1.2em}{\includegraphics{images/antiquark-flowline.pdf}}
  & = &
  \begin{split}
    \delta_{\alpha\beta} \delta_{ij} \, \theta(t-s) \, \widetilde K(t-s,p) \,.
  \end{split}
  & &
  \label{eq:lambdabarchi}
\end{align}
Since $\widetilde K(t,p)$ only depends on $p$ only quadratically, the
momentum direction for the fermion flow lines is irrelevant.


\subsubsection*{Flow vertices}

Since $\mathcal{L}_B$ and $\mathcal{L}_\chi$ of
\eqn{eq:Lagrangian_flow} are proportional to a Lagrange multiplier
field, the resulting vertices always involve at least one flow
line. Such vertices are always associated with a flow-time
parameter which is integrated over. We denote them by ``flow vertices''
and represent them by empty circles in Feynman diagrams. The
corresponding Feynman rules can be derived straightforward.
In this paper, we define Feynman rules by assuming all momenta to be
outgoing.

The three-point gluon flow vertex is governed by $X_2$:
\begin{align}
  &
  \begin{gathered}
    \includegraphics{images/ggg-flow-vertex.pdf}
  \end{gathered}
  & = &
  \begin{split}
    -\mathrm{i} g f^{abc} \int_0^\infty \mathrm{d}s \, \big( & \delta_{\nu\rho} (r-q)_\mu + 2 \delta_{\mu\nu} q_\rho - 2 \delta_{\mu\rho} r_\nu \\
    & + (\kappa-1) (\delta_{\mu\rho} q_\nu  - \delta_{\mu\nu} r_\rho) \big) \,. \\
  \end{split}
  & &
  \label{eq:3gvert}
\end{align}
The interaction terms of $\mathcal{L}_B$ involve exactly one Lagrange
multiplier field $L(s)$, cf.\ \eqn{eq:Lagrangian_flow}.  According to
\eqs{eq:bl} and \noeqn{eq:blfr}, it must be contracted with a flowed
gluon field $B(t)$ at larger flow time $t>s$. Thus, the vertex in
\eqn{eq:3gvert} contains exactly one outgoing flow line. On the other
hand, each of the two flowed gluon fields $B(s)$ in the interaction
terms of $\mathcal{L}_B$ can be contracted either with a Lagrange
multiplier field $L(t')$ at \textit{smaller} flow time $t'<s$, resulting
in an ingoing flow line, or with another flowed gluon field $B(s')$. In
this sense, the Feynman rule \noeqn{eq:3gvert} actually represents three
vertices, displayed in \fig{fig:3gcomb}: They all contain one outgoing
flow line, while each of the other two lines can either be an ingoing
flow line or a flowed gluon propagator.\footnote{We focus on the
  calculation of Green's functions in this paper; if one calculates
  amplitudes, the lines could also represent external ``particles'', of
  course.}  The dashed arrows in \eqn{eq:3gvert} indicate lines which
can be both a flow line or a propagator, while the solid arrows always
denote flow lines.


\begin{figure}
  \begin{equation*}
    \begin{gathered}
      \includegraphics{images/ggg-flow-vertex_fff.pdf}
    \end{gathered}
    \qquad
    \begin{gathered}
      \includegraphics{images/ggg-flow-vertex_gff.pdf}
    \end{gathered}
    \qquad
    \begin{gathered}
      \includegraphics{images/ggg-flow-vertex_ggf.pdf}
    \end{gathered}
  \end{equation*}
  \caption{The three versions of the vertex $X_2$. Lines with an
    adjacent arrow denote flow lines, the others are flowed (or regular)
    gluons.}
  \label{fig:3gcomb}
\end{figure}


Similarly, the four-point gluon flow vertex is governed by $X_3$:
\begin{align}
  &
  \begin{gathered}
    \includegraphics{images/gggg-flow-vertex.pdf}
  \end{gathered}
  & = &
  \begin{split}
    -g^2 \int_0^\infty \mathrm{d}s \, \big( & f^{abe} f^{cde} (\delta_{\mu\rho}\delta_{\nu\sigma} - \delta_{\mu\sigma}\delta_{\nu\rho} ) \\
    & + f^{ace} f^{bde} (\delta_{\mu\nu}\delta_{\rho\sigma} - \delta_{\mu\sigma}\delta_{\nu\rho} ) \\
    & + f^{ade} f^{bce} (\delta_{\mu\nu}\delta_{\rho\sigma} - \delta_{\mu\rho}\delta_{\nu\sigma} ) \big) \,.
  \end{split}
  & &
  \label{eq:4gvert}
\end{align}
Again, there is exactly one outgoing flow line, while the other three
lines are either flowed gluons or incoming flow lines. This means that
the Feynman rule \noeqn{eq:4gvert} actually represents four different
vertices.

Note that the factor $2\TR$, arising from the trace in
\eqn{eq:Lagrangian_flow}, has been discarded in both vertices in accordance with the
normalization of \eqn{eq:blfr}.

Similar considerations applied to $\mathcal{L}_\chi$ lead to vertices
involving flowed quark fields. For example, the quark flow vertex with
one gluon is described by the following Feynman rule:
\begin{align}
  &
  \begin{gathered}
    \includegraphics{images/qqg-flow-vertex.pdf}
  \end{gathered}
  & = &
  \begin{split}
    \mathrm{i} g \, \delta_{\alpha\beta} \big( T^a \big)_{ij} \int_0^\infty \mathrm{d} s \, \big( 2 p_\mu + (1-\kappa) q_\mu \big) \,.
  \end{split}
  & &
\end{align}
In this case, the fermion line with outgoing charge flow is also
outgoing in the gradient flow, while the other two lines can be either
flowed propagators, or ingoing flow lines. Thus, this Feynman rule
actually represents four different vertices.

A complete list of the Feynman rules can be found in
Appendix\,\ref{sec:frules}.

Let us summarize the Feynman rules for the \gff\ and point out a few
more features:
\begin{enumerate}
  \item Propagators and flow lines always carry an exponential factor
    $\mathrm{e}^{\pm tp^2}$ for each vertex, where $t$ is the flow time
    of the vertex.
  \item Flow lines always start at a flow vertex.  They end either at
    another flow vertex or an external operator. The $\theta$-distribution
    $\theta (t-s)$ implies a flow-time direction, which we denote by an
    adjacent arrow in Feynman diagrams.
  \item A flow vertex always has an outgoing flow line connected to it.
    The other lines are either incoming flow lines, or flowed
    propagators.  Each flow vertex is defined at a
    flow time $s$ and implies an integration $\int_0^\infty
    \mathrm{d}s$.  The $\theta$-distribution of the outgoing flow line
    restricts the integration to a finite upper limit $t$,
    i.e.\ $\int_0^t \mathrm{d}s$.
  \item Diagrams with closed flow-line loops vanish, because the
    integration interval shrinks to a point. Hence, only diagrams whose
    flow lines form trees contribute to any observable.
\end{enumerate}
